	\documentclass[14pt]{extreport}
\usepackage{extsizes}
	\usepackage[frenchb]{babel}
	\usepackage[utf8]{inputenc}  
	\usepackage[T1]{fontenc}
	\usepackage{amssymb}
	\usepackage[mathscr]{euscript}
	\usepackage{stmaryrd}
	\usepackage{amsmath}
	\usepackage{tikz}
	\usepackage[all,cmtip]{xy}
	\usepackage{amsthm}
	\usepackage{varioref}
	\usepackage[ margin=1in]{geometry}
	\geometry{a4paper}
	\usepackage{lmodern}
	\usepackage{hyperref}
	\usepackage{array}
	\usepackage{float}
	\usepackage{easytable}
\renewcommand{\theenumi}{\alph{enumi})}
	 \usepackage{fancyhdr}\usepackage{longtable}
	 \usetikzlibrary{shapes.misc}

\tikzset{cross/.style={cross out, draw=black, minimum size=2*(#1-\pgflinewidth), inner sep=0pt, outer sep=0pt},
%default radius will be 1pt. 
cross/.default={1pt}}

	\pagestyle{fancy}
	\theoremstyle{plain}
	\fancyfoot[C]{\empty} 
	\fancyhead[L]{Contrôle}
	\fancyhead[R]{10 mars 2026}
	
	
	\title{Contrôle chapitre 5}
	\date{}
	\begin{document}

\begin{center}{\Large Contrôle chapitre 5}\end{center}

 
 
\subsection*{Exercice 1 (5 points)}
 
 
 Trier dans l'ordre croissant les nombres suivants : 
 \[ A = 3 - (-5) ;\phantom{miaou}
  B = 5 - (-2) ;\phantom{miaou}
   C = 3 - 5 ; \phantom{miaou}
   D = 5 -3  ;\phantom{miaou}
   E = -5 - 3\]
\subsection*{Exercice 2 (5 points)}
 
 Calculez : \begin{enumerate}
 \item $43 + (-12) + (-14) + (+7) $
 \item $-3 + 2 - (-1)$
 \item $ 52 - (34 - (-2)) + (-12 + 3)$.
 \item $ [ 8 - (3 - 5)] + [(12-14) -3]$.
 \item $ (-1,5) + (3,4)  + (-11,7) + (-1,3)$ 
 \end{enumerate}

    
    \subsection*{Exercice 3 (4 points) }
    
Compléter le carré magique suivant avec des entiers relatifs pour que les quatre lignes, les quatre colonnes, 
et les deux diagonales aient une somme égale à $6$. 
\[
\begin{TAB}(e,1cm){|c|c|c|c|}{|c|c|c|c|}
    1 & & &  \\
    1 &  & -3 &  \\
      & 2 & 0 & \\
    5  & -7 &  & -2
\end{TAB}
\]


\subsection*{Exercice 4 (6 points)}

\begin{figure}[H]\center
\begin{tikzpicture}
\draw[dotted](-5.5,-4.5) grid (7.5,6.5);
\draw[very thick, ->] (-5.5, 0)--(7.5,0);
\draw[very thick, ->] (0,-4.5)--(0,6.5);
\fill (1,-3) node[cross=5pt, rotate = 0] {};
\draw (1,-3) ++ (.25, .25) node {$A$};
\fill (-5,1) node[cross=5pt, rotate = 0] {};
\draw (-5,1) ++ (.25, .25) node {$B$};
\fill (4,5) node[cross=5pt, rotate = 0] {};
\draw (4,5) ++ (.25, .25) node {$C$};
\draw (1,.3) node {$1$};
\draw (.3,1) node {$1$};
\end{tikzpicture}
\end{figure}
1)[1,5 pts] Indiquer les coordonnées des points $A$, $B$ et $C$. 

2)[1 pt] Placer le point $D$ de coordonnées $(2; 5)$ et le point $E$ de coordonnées $(-3; 2)$. 

3)[1 pt] Placer le point $F$ qui a la même abscisse que $A$ et la même 
ordonnée que $B$. 



4)[1 pt] Placer le point $G$ dont l'abscisse est le tiers de la somme des 
abscisses de $A$, $B$ et $C$, et dont l'ordonnée est le tiers de la 
somme des ordonnées de $A$, $B$ et $C$. 

5)[1 pt] Tracer les droites $(AG)$, $(BG)$ et $(CG)$, et le triangle $ABC$. Que représentent les droites ainsi tracées pour le triangle $ABC$ ?
\newpage 

\begin{center}{\Large Contrôle chapitre 5}\end{center}

 
 
\subsection*{Exercice 1 (5 points)}
 
 
 Trier dans l'ordre croissant les nombres suivants : 
 \[ A = 4 - (-5) ;\phantom{miaou}
  B = 2 - (-2) ;\phantom{miaou}
   C = 3 - 5 ; \phantom{miaou}
   D = 5 -3  ;\phantom{miaou}
   E = -5 - 3\]
\subsection*{Exercice 2 (5 points)}
 
 Calculez : \begin{enumerate}
 \item $33 + (-12) + (-14) + (+7) $
 \item $-5 + 2 - (-1)$
 \item $ 32 - (34 - (-2)) + (-12 + 3)$.
 \item $ [ 18 - (3 - 5)] + [(12-14) -3]$.
 \item $ (-2,5) + (3,4)  + (-11,7) + (-1,3)$ 
 \end{enumerate}

    
    \subsection*{Exercice 3 (4 points) }
    
Compléter le carré magique suivant avec des entiers relatifs pour que les quatre lignes, les quatre colonnes, 
et les deux diagonales aient une somme égale à $6$. 
\[
\begin{TAB}(e,1cm){|c|c|c|c|}{|c|c|c|c|}
    5 & & &  \\
    -1 &  & 0 &  \\
      & 7 & -3 & \\
    1  & 4 &  & 2
\end{TAB}
\]


\subsection*{Exercice 4 (6 points)}

\begin{figure}[H]\center
\begin{tikzpicture}
\draw[dotted](-5.5,-4.5) grid (7.5,6.5);
\draw[very thick, ->] (-5.5, 0)--(7.5,0);
\draw[very thick, ->] (0,-4.5)--(0,6.5);
\fill (1,-3) node[cross=5pt, rotate = 0] {};
\draw (1,-3) ++ (.25, .25) node {$A$};
\fill (-5,4) node[cross=5pt, rotate = 0] {};
\draw (-5,4) ++ (.25, .25) node {$B$};
\fill (4,2) node[cross=5pt, rotate = 0] {};
\draw (4,2) ++ (.25, .25) node {$C$};
\draw (1,.3) node {$1$};
\draw (.3,1) node {$1$};
\end{tikzpicture}
\end{figure}
1)[1,5 pts] Indiquer les coordonnées des points $A$, $B$ et $C$. 

2)[1 pt] Placer le point $D$ de coordonnées $(3; 6)$ et le point $E$ de coordonnées $(-3; -2)$. 

3)[1 pt] Placer le point $F$ qui a la même abscisse que $A$ et la même 
ordonnée que $B$. 



4)[1 pt] Placer le point $G$ dont l'abscisse est le tiers de la somme des 
abscisses de $A$, $B$ et $C$, et dont l'ordonnée est le tiers de la 
somme des ordonnées de $A$, $B$ et $C$. 

5)[1 pt] Tracer les droites $(AG)$, $(BG)$ et $(CG)$, et le triangle $ABC$. Que représentent les droites ainsi tracées pour le triangle $ABC$ ?
\newpage 

 
\end{document}